% !TITLE 黄鹤楼送孟浩然之广陵
% !AUTHOR 李白
% !DYNASTY 唐
% !ORDER 7

\begin{poem}
故人西辞黄鹤楼\textsuperscript{1}，烟花三月\textsuperscript{2}下扬州\textsuperscript{3}。
孤帆远影碧空尽\textsuperscript{4}，唯见长江天际流。
\end{poem}

\begin{annotations}
\notetext{1}{黄鹤楼：在今天的武汉，长江边上的名楼，古代文人常在此聚会送别。}
\notetext{2}{烟花三月：春天三月，江南水汽朦胧、花开繁盛。“烟花”不是烟火，而是指烟雾朦胧中的花景。}
\notetext{3}{扬州：今江苏扬州。唐朝时扬州是天下最繁华的商业都市之一，有“扬一益二”之说（扬州第一，益州第二）。孟浩然从武汉乘船顺长江东下到扬州。}
\notetext{4}{碧空尽：船帆的影子在碧蓝的天空尽头消失了。尽，指完全消失在天际线的尽头。}
\end{annotations}

\begin{appreciation}
《黄鹤楼送孟浩然之广陵》是李白最出名的送别诗之一，也大概是他写得最不使劲的一首。孟浩然比李白大十二岁，是李白佩服的前辈和朋友。广陵就是扬州，在长江下游。李白在黄鹤楼（今天的武汉）送孟浩然上船，顺流东下。

“故人西辞黄鹤楼，烟花三月下扬州。”从西边的黄鹤楼辞别，在三月里顺流而下去扬州。烟花不是烟火，是江南三月的水汽与繁花混在一起、朦胧胧的样子。扬州是当时天下最繁华的城市之一。李白把朋友送去了一个好地方，还是最好的季节——这首诗一开始就没有眼泪。

“孤帆远影碧空尽，唯见长江天际流。”孤帆，不是江上只有一艘船，是李白眼里只有这一艘船。船越走越远，帆影在碧蓝的天空尽头消失；看不见船了，只剩长江往天际流去。“尽”字写得精确——不是模糊地看不见了，而是消失在一个可以凝视的位置上。然后是“唯见”：李白还站在那里。

这首诗从头到尾没写一个“舍不得”，全是“看”。可正是这个一直看到看不见的人，把什么都说尽了。还记得《行行重行行》吗？那里头“行行重行行”，是游子越走越远，家中的女子盼他回头；《蒹葭》里的人顺着水去追，伊人“宛在水中央”，追不上。这里是反过来的：一个人看着朋友越走越远，不追，就送到底。古诗十九首里的人等的是回来，李白送的是远去——同样是“越走越远”，一盼一送，都是中国文学里最深的情意。

将来你去很远的地方上学，我大概也会在机场站到看不见为止。你不用回头，往前走就行。
\end{appreciation}
